\scanpage{234}
\mspage{- 189 -}\hfill IX

Donc les structures d'extension (14.1.2) et (14.1.1) peuvent s'interpréter respectivement (VIII 3.7) comme équivalant à la donnée d'homomorphisme
\[
\underline M_K\xrightarrow{p_K} B_K\ ,\quad \underline M'_K\xrightarrow{p'_K} B'_K\ .\tag*{(14.1.4)}
\]

Soient $W_K^{\handwritten{\unsure{\natural|\leftrightarrows}}}$ la biextension de $(B_K,B'_K)$ par $\mathbb G_{mK}$ qui définit la dualité canonique entre ces schémas abéliens (8.2), et $E_K$ la biextension de $(\underline M_K,\underline M'_K)$ par $\mathbb G_{mK}$ image inverse de $W_K^{\handwritten{\natural}}$ par les homomorphismes (14.1.4). Notons qu'on vérifie immédiatement que les trivialisations de $E_K$ correspondent biunivoquement, de façon canonique, aux homomorphismes
\[
i:\underline M_K\longrightarrow G_K\tag*{(14.1.5)}
\]
qui relèvent $p_K:\underline M_K\longrightarrow B_K$, ou, symétriquement, aux homomorphismes
\[
i':\underline M'_K\longrightarrow G'_K\tag*{(14.1.6)}
\]
qui relèvent $P'_K:\underline M'_K\longrightarrow B'_K$.

Ceci posé, désignant par un exposant $^{\mathrm{an}}$ les groupes rigide-analytiques sur $K$ associés aux groupes algébriques envisagés [13], on définit dans [24] une suite exacte canonique
\[
0\longrightarrow \underline M_K^{\mathrm{an}}\xrightarrow{i^{\mathrm{an}}}G_K^{\mathrm{an}}\xrightarrow{e} A_K^{\mathrm{an}}\longrightarrow 0\ ,\tag*{(14.1.7)}
\]
dans laquelle l'homomorphisme $i^{\mathrm{an}}$, qui peut aussi s'interpréter comme un homomorphisme $i:\underline M_K\longrightarrow G_K$ pour des schémas en groupes ordinaires, \uline{relève l'homomorphisme} $p$ (14.1.4), i.e. est de la forme envisagée dans (14.1.5). Cette suite exacte est l'analogue rigide-analytique de la suite exacte (13.3.4). Elle dépend de façon fonctorielle de $A_K$, et

\handwritten{507}

\mspage{507}
% archive sheet number: 233

